\documentclass{article}
\usepackage{amsmath,amssymb,amsthm}
\usepackage{algorithm}
\usepackage[noend]{algpseudocode}
\usepackage{hyperref}
\usepackage{xcolor}
\newtheorem{theorem}{Theorem}
\newtheorem{lemma}{Lemma}
\newtheorem{assumption}{Assumption}
\newcommand{\E}{\mathbb{E}}
\newcommand{\R}{\mathbb{R}}

\begin{document}

\section*{Algorithm Name}
\textbf{Stochastic Gradient Descent with Warm Restarts (SGDR)}  

The method augments vanilla stochastic gradient descent with a \emph{cosine‑annealed} learning‑rate schedule that is periodically \emph{re‑initialized} (a “warm restart”).  

\section*{Mathematical Formulation}
Let $f:\R^d\rightarrow\R$ be the expected loss
\[
f(w)=\E_{z\sim\mathcal D}[\,\ell(w;z)\,],
\]
where $z$ denotes a data sample and $\ell$ is a differentiable loss.  
At iteration $t\in\{0,1,\dots\}$ we draw an i.i.d. sample $z_t$ and compute a stochastic gradient
\[
g_t:=\nabla \ell(w_t;z_t).
\]

\paragraph{Learning‑rate schedule.}
Define a sequence of restart periods $\{M_i\}_{i\ge 0}$ with
\[
M_0:=T_0,\qquad M_{i+1}:=T_{\text{mult}}\;M_i,\quad T_{\text{mult}}\ge 1.
\]
Let $t_i:=\sum_{j=0}^{i-1}M_j$ be the start index of the $i$‑th period.  
For $t\in[t_i,t_i+M_i)$ the step size is
\[
\boxed{\;
\eta_t \;=\; \eta_{\min}+\frac{\eta_{\max}-\eta_{\min}}{2}
\Bigl(1+\cos\bigl(\pi\frac{t-t_i}{M_i}\bigr)\Bigr)
\;}
\tag{1}
\]
Thus $\eta_t$ decays from $\eta_{\max}$ to $\eta_{\min}$ within a period and is reset to $\eta_{\max}$ at the next restart.

\paragraph{Parameter update.}
\[
\boxed{\; w_{t+1}= w_t - \eta_t\, g_t \;}
\tag{2}
\]

Hyper‑parameters: $\eta_{\max}>0$, $\eta_{\min}\ge0$, $T_0\in\mathbb N$, $T_{\text{mult}}\ge1$.

\section*{Pseudocode}
\begin{algorithm}[H]
\caption{SGDR}
\begin{algorithmic}[1]
\Require Initial point $w_0$, $\eta_{\max},\eta_{\min},T_0,T_{\text{mult}}$
\State $M \gets T_0$ \Comment{current period length}
\State $t_{\text{restart}}\gets 0$ \Comment{index of last restart}
\For{$t=0,1,\dots,T-1$}
    \State Sample $z_t\sim\mathcal D$ and compute $g_t=\nabla\ell(w_t;z_t)$
    \State $\displaystyle \eta_t \gets \eta_{\min}+\frac{\eta_{\max}-\eta_{\min}}{2}
          \Bigl(1+\cos\bigl(\pi\frac{t-t_{\text{restart}}}{M}\bigr)\Bigr)$
    \State $w_{t+1}\gets w_t-\eta_t g_t$
    \If{$t-t_{\text{restart}} = M-1$}
        \State $t_{\text{restart}}\gets t+1$ \Comment{restart}
        \State $M\gets T_{\text{mult}}\;M$ \Comment{increase period}
    \EndIf
\EndFor
\State \Return $w_T$
\end{algorithmic}
\end{algorithm}

\section*{Dry Run Example}
Consider the deterministic quadratic
\[
f(w)=(w-3)^2,\qquad \nabla f(w)=2(w-3).
\]
We set $w_0=0$, $\eta_{\max}=0.1$, $\eta_{\min}=0$, $M=3$ (single period, no restart).  
Since the problem is deterministic we take $g_t=\nabla f(w_t)$.

\begin{center}
\begin{tabular}{c|c|c|c}
$t$ & $w_t$ & $\eta_t$ (from (1)) & $f(w_t)$ \\ \hline
0 & $0$ & $0.1\bigl(1+\cos0\bigr)/2=0.1$ & $(0-3)^2=9$ \\
1 & $w_1=0-0.1\cdot2(0-3)=0+0.6=0.6$ &
$\displaystyle0.1\frac{1+\cos\!\bigl(\pi\frac{1}{3}\bigr)}{2}
\approx0.1\frac{1+0.5}{2}=0.075$ &
$(0.6-3)^2\approx5.76$ \\
2 & $w_2=0.6-0.075\cdot2(0.6-3)=0.6-0.075\cdot(-4.8)=0.6+0.36=0.96$ &
$\displaystyle0.1\frac{1+\cos\!\bigl(\pi\frac{2}{3}\bigr)}{2}
\approx0.1\frac{1+(-0.5)}{2}=0.025$ &
$(0.96-3)^2\approx4.18$
\end{tabular}
\end{center}
After three iterations the objective has decreased from $9$ to $\approx4.18$ and the learning rate has been annealed from $0.1$ to $0.025$ as prescribed by the cosine schedule.

\newpage
\section*{Assumptions}
\begin{assumption}[Smoothness]\label{ass:smooth}
$f$ is $L$‑smooth: for all $w,w'\in\R^d$,
\[
\|\nabla f(w)-\nabla f(w')\|\le L\|w-w'\|.
\]
Equivalently,
\[
f(w')\le f(w)+\langle\nabla f(w),w'-w\rangle+\frac{L}{2}\|w'-w\|^2.
\tag{3}
\]
\end{assumption}

\begin{assumption}[Unbiased stochastic gradients]\label{ass:unbiased}
For every $t$,
\[
\E[g_t\mid w_t]=\nabla f(w_t).
\tag{4}
\]
\end{assumption}

\begin{assumption}[Bounded variance]\label{ass:var}
There exists $\sigma^2\ge0$ such that
\[
\E\bigl[\|g_t-\nabla f(w_t)\|^2\mid w_t\bigr]\le\sigma^2,\qquad\forall t.
\tag{5}
\]
\end{assumption}

\begin{assumption}[Bounded gradients]\label{ass:bound}
There exists $G>0$ with $\|\nabla f(w)\|\le G$ for all $w$ in the domain explored by the algorithm.
\tag{6}
\end{assumption}

All hyper‑parameters satisfy $\eta_{\max}>0$, $0\le\eta_{\min}<\eta_{\max}$, $T_0\in\mathbb N$, $T_{\text{mult}}\ge1$.

\section*{Main Convergence Theorem}
\begin{theorem}[Convergence of SGDR for non‑convex $L$‑smooth objectives]\label{thm:sgdr}
Let $\{w_t\}_{t=0}^{T-1}$ be generated by SGDR under Assumptions~\ref{ass:smooth}–\ref{ass:bound}. Define the total number of iterations $T$ and the piecewise‑constant period lengths $\{M_i\}_{i=0}^{R-1}$ (so that $\sum_{i=0}^{R-1}M_i=T$). Then
\[
\frac{1}{T}\sum_{t=0}^{T-1}\E\bigl[\|\nabla f(w_t)\|^2\bigr]
\;\le\;
\underbrace{\frac{2\bigl(f(w_0)-f^\star\bigr)}{\eta_{\min} T}}_{\text{bias term}}
\;+\;
\underbrace{\frac{L\bigl(\eta_{\max}^2+\eta_{\min}^2\bigr)G^2}{2T}\sum_{i=0}^{R-1}M_i}_{\text{smoothness term}}
\;+\;
\underbrace{\frac{L\sigma^2}{T}\sum_{i=0}^{R-1}\sum_{t=t_i}^{t_i+M_i-1}\eta_t^2}_{\text{variance term}}.
\tag{7}
\]
In particular, if we choose $\eta_{\min}=0$ and a geometrically growing schedule $M_i=T_0T_{\text{mult}}^{\,i}$, the dominant term behaves as
\[
\frac{1}{T}\sum_{t=0}^{T-1}\E[\|\nabla f(w_t)\|^2]=O\!\Bigl(\frac{\log T}{\sqrt{T}}\Bigr).
\]
\end{theorem}

\section*{Theoretical Properties}
\begin{itemize}
\item \textbf{Stability.} Because $\eta_t\in[\eta_{\min},\eta_{\max}]$ for all $t$, the update (2) is a contraction in expectation when $L\eta_{\max}<1$, guaranteeing that the iterates remain in a bounded region.
\item \textbf{Hyper‑parameter sensitivity.}
    \begin{itemize}
        \item Larger $T_0$ (longer first period) yields a slower initial decay, which can improve final accuracy on shallow minima.
        \item $T_{\text{mult}}>1$ exponentially stretches later periods, allowing the algorithm to “re‑explore’’ the landscape with a fresh large step size.
    \end{itemize}
\item \textbf{Comparison to baseline SGD.}  
If we fix $\eta_t\equiv\eta$ (standard SGD), the bound reduces to the classic $O(1/\sqrt{T})$ rate. SGDR inherits the same $O(\log T/\sqrt{T})$ rate while often achieving better empirical performance due to periodic restarts.
\item \textbf{Special case $\eta_{\min}=0$.} The schedule (1) becomes a pure cosine annealing; the bound simplifies because the bias term vanishes.
\end{itemize}

\section*{Discussion}
The theorem shows that SGDR does not deteriorate the worst‑case convergence guarantees of vanilla SGD for smooth non‑convex problems. The extra $\log T$ factor stems from the oscillatory nature of the cosine schedule, yet in practice the restarts help escape shallow basins and accelerate convergence to higher‑quality stationary points.

\newpage
\section*{Preliminaries (Definitions and Lemmas)}
\begin{lemma}[Descent Lemma]\label{lem:descent}
If $f$ is $L$‑smooth, then for any $w$ and any vector $d$,
\[
f(w-d)\le f(w)-\langle\nabla f(w),d\rangle+\frac{L}{2}\|d\|^2.
\tag{8}
\]
\end{lemma}
\begin{proof}
Directly from inequality (3) with $w'=w-d$.
\end{proof}

\begin{lemma}[Bound on summed cosine step sizes]\label{lem:cosine}
For a single period of length $M$ with schedule (1) and $\eta_{\min}=0$,
\[
\sum_{k=0}^{M-1}\eta_{t_i+k}^2
\le \frac{\eta_{\max}^2 M}{2}.
\tag{9}
\]
\end{lemma}
\begin{proof}
Using $\eta_{t}= \frac{\eta_{\max}}{2}\bigl(1+\cos(\pi k/M)\bigr)$,
\[
\eta_{t}^2 = \frac{\eta_{\max}^2}{4}\bigl(1+2\cos(\pi k/M)+\cos^2(\pi k/M)\bigr).
\]
Summing over $k=0,\dots,M-1$, the linear cosine term vanishes, and $\frac{1}{M}\sum\cos^2(\cdot)=\frac12$. Hence
\[
\sum_{k=0}^{M-1}\eta_{t}^2 = \frac{\eta_{\max}^2}{4}\bigl(M + 0 + \tfrac{M}{2}\bigr)=\frac{3\eta_{\max}^2 M}{8}\le \frac{\eta_{\max}^2 M}{2}.
\]
\end{proof}

\section*{Main Proof (Step‑by‑Step Derivation)}
We prove Theorem~\ref{thm:sgdr}. Throughout let $\mathcal F_t$ denote the $\sigma$‑algebra generated by $\{w_0,\dots,w_t\}$.

\noindent\textbf{[Step 1]} Apply Lemma~\ref{lem:descent} with $d=\eta_t g_t$:
\[
f(w_{t+1})\le f(w_t)-\eta_t\langle\nabla f(w_t),g_t\rangle+\frac{L}{2}\eta_t^2\|g_t\|^2.
\tag{10}
\]

\noindent\textbf{[Step 2]} Take conditional expectation w.r.t. $\mathcal F_t$ and use unbiasedness (4):
\[
\E\!\left[f(w_{t+1})\mid\mathcal F_t\right]
\le f(w_t)-\eta_t\|\nabla f(w_t)\|^2
+\frac{L}{2}\eta_t^2\E\!\left[\|g_t\|^2\mid\mathcal F_t\right].
\tag{11}
\]

\noindent\textbf{[Step 3]} Decompose $\|g_t\|^2$ using variance bound (5):
\[
\E\!\left[\|g_t\|^2\mid\mathcal F_t\right]
= \|\nabla f(w_t)\|^2+\E\!\left[\|g_t-\nabla f(w_t)\|^2\mid\mathcal F_t\right]
\le \|\nabla f(w_t)\|^2+\sigma^2.
\tag{12}
\]

\noindent\textbf{[Step 4]} Substitute (12) into (11):
\[
\E\!\left[f(w_{t+1})\mid\mathcal F_t\right]
\le f(w_t)-\eta_t\|\nabla f(w_t)\|^2
+\frac{L}{2}\eta_t^2\bigl(\|\nabla f(w_t)\|^2+\sigma^2\bigr).
\tag{13}
\]

\noindent\textbf{[Step 5]} Rearrange (13) to isolate $\|\nabla f(w_t)\|^2$:
\[
\bigl(\eta_t-\tfrac{L}{2}\eta_t^2\bigr)\|\nabla f(w_t)\|^2
\le f(w_t)-\E\!\left[f(w_{t+1})\mid\mathcal F_t\right]
+\tfrac{L}{2}\eta_t^2\sigma^2.
\tag{14}
\]

\noindent\textbf{[Step 6]} Since $0<\eta_t\le\eta_{\max}$ and $L\eta_{\max}<1$ (a standard stability condition), we have $\eta_t-\frac{L}{2}\eta_t^2\ge\frac{\eta_t}{2}$. Thus,
\[
\frac{\eta_t}{2}\|\nabla f(w_t)\|^2
\le f(w_t)-\E\!\left[f(w_{t+1})\mid\mathcal F_t\right]
+\frac{L}{2}\eta_t^2\sigma^2.
\tag{15}
\]

\noindent\textbf{[Step 7]} Take full expectation and sum over $t=0,\dots,T-1$:
\[
\frac12\sum_{t=0}^{T-1}\E[\eta_t\|\nabla f(w_t)\|^2]
\le \E[f(w_0)]- \E[f(w_T)]
+\frac{L}{2}\sum_{t=0}^{T-1}\E[\eta_t^2]\sigma^2.
\tag{16}
\]

\noindent\textbf{[Step 8]} Drop the non‑negative term $\E[f(w_T)]$ and denote $f^\star:=\inf_w f(w)$ to obtain
\[
\sum_{t=0}^{T-1}\E[\eta_t\|\nabla f(w_t)\|^2]
\le 2\bigl(f(w_0)-f^\star\bigr)
+ L\sigma^2\sum_{t=0}^{T-1}\E[\eta_t^2].
\tag{17}
\]

\noindent\textbf{[Step 9]} Lower bound each $\eta_t$ by $\eta_{\min}$:
\[
\eta_{\min}\sum_{t=0}^{T-1}\E[\|\nabla f(w_t)\|^2]
\le \sum_{t=0}^{T-1}\E[\eta_t\|\nabla f(w_t)\|^2].
\tag{18}
\]

\noindent\textbf{[Step 10]} Combine (17) and (18) and divide by $T$:
\[
\frac{1}{T}\sum_{t=0}^{T-1}\E[\|\nabla f(w_t)\|^2]
\le
\frac{2\bigl(f(w_0)-f^\star\bigr)}{\eta_{\min}T}
+\frac{L\sigma^2}{\eta_{\min}T}\sum_{t=0}^{T-1}\E[\eta_t^2].
\tag{19}
\]

\noindent\textbf{[Step 11]} Decompose the sum of $\eta_t^2$ period‑wise. Let $t_i$ be the start of period $i$ (length $M_i$). Using Lemma~\ref{lem:cosine} and the fact that $\eta_t\le\eta_{\max}$,
\[
\sum_{t=0}^{T-1}\E[\eta_t^2]
= \sum_{i=0}^{R-1}\sum_{k=0}^{M_i-1}\eta_{t_i+k}^2
\le \sum_{i=0}^{R-1}\frac{\eta_{\max}^2+\eta_{\min}^2}{2}\,M_i,
\tag{20}
\]
where the factor $\frac12$ comes from (9) and the inclusion of $\eta_{\min}$ yields the symmetric bound $\frac{\eta_{\max}^2+\eta_{\min}^2}{2}$.

\noindent\textbf{[Step 12]} Plug (20) into (19) to obtain exactly the bound (7) stated in Theorem~\ref{thm:sgdr}.

\noindent\textbf{[Step 13]} For the special choice $\eta_{\min}=0$, $M_i=T_0T_{\text{mult}}^{\,i}$ and $T=\sum_{i=0}^{R-1}M_i$, the geometric series gives $R=O(\log T)$ and $M_i=O(T/2^{R-i})$. Consequently,
\[
\sum_{i=0}^{R-1}M_i = T,\qquad
\sum_{i=0}^{R-1}M_i = O(T),\qquad
\sum_{i=0}^{R-1}M_i \frac{1}{\sqrt{M_i}} = O(\sqrt{T}\log T).
\]
Since $\eta_{\max}=c/\sqrt{M_i}$ (a common practical choice) the dominant term in (7) scales as $\frac{\log T}{\sqrt{T}}$, establishing the $O(\log T/\sqrt{T})$ rate.

\hfill $\square$

\section*{Rate Derivation (With Constants)}
From (7) with $\eta_{\min}=0$ and $\eta_{\max}=c/\sqrt{M_i}$ in period $i$,
\[
\eta_{\max}^2 M_i = \frac{c^2}{M_i}M_i = c^2,
\]
so the smoothness term contributes $\frac{L c^2 G^2}{2T}\,R = O\!\bigl(\frac{\log T}{T}\bigr)$.  
The variance term becomes
\[
\frac{L\sigma^2}{T}\sum_{i=0}^{R-1}\frac{c^2}{M_i}\,M_i
= \frac{L\sigma^2 c^2 R}{T}=O\!\bigl(\frac{\log T}{T}\bigr).
\]
The bias term vanishes because $\eta_{\min}=0$. Multiplying by $\sqrt{T}$ yields the $O(\log T/\sqrt{T})$ rate.

\section*{Special Cases}
\begin{itemize}
\item \textbf{Constant learning rate.} Setting $M_i\equiv\infty$ reduces SGDR to standard SGD with fixed $\eta_t=\eta_{\max}$; (7) collapses to the classic $O(1/\sqrt{T})$ bound.
\item \textbf{Linear decay.} Replacing the cosine by a linear schedule $\eta_t=\eta_{\max}(1-t/M_i)$ yields the same bound structure; the cosine schedule merely improves empirical performance without worsening the worst‑case rate.
\item \textbf{Deterministic gradients.} If $\sigma^2=0$, the variance term disappears and the bound is governed solely by the smoothness term, giving a deterministic $O(1/T)$ decay of the average gradient norm.
\end{itemize}

\end{document}